\section{Introduction}
Two PDEs can look almost the same on paper and still produce very different solutions. The reverse is also common: one problem can be written in many equivalent ways. A representation that helps scientists find related problems should cope with both situations. It should place problems with similar solutions close together, and it should not be thrown off by a change of notation.

General-purpose embedding models are already used to search and organize scientific text and images, and many of them accept text, images, or both. It is natural to ask whether their similarity scores say anything about how PDE solutions behave. We ask a narrow version of this question: \emph{when an embedding model rates one pair of problems as more similar than another, do the solutions of the first pair actually evolve more similarly?} A positive answer on our benchmark would not show general mathematical reasoning. A negative answer would show that these embedding spaces should not be read as statements about dynamics.

Answering the question requires a ground truth, which we build by simulation. We take six equations with two fields on the unit square, three families of initial conditions (ICs), and three boundary conditions (BCs). We call each combination of equation, IC family, and BC a \emph{problem}. For each of the 54 problems we simulate 100 samples, giving 5,400 trajectories. Problems in the same IC family start from exactly the same initial fields. This lets us compare them directly: start two systems from the same state and measure how far apart their solutions drift. Our reference distance is the root-mean-square (RMS) difference between the fields.

We then query three multimodal embedding models, used as released and without fine-tuning, in three ways: with a text description of the problem, with an image of five solution snapshots, and with both together. In each case we compare the models' similarity rankings with the reference rankings. Simple baselines and a set of controls help interpret the results.

The answer is mostly negative. Images give moderate agreement with the reference, but comparing the raw pixels gives more. Text descriptions give weak agreement, about the same as matching character sequences. They are also fragile: for two of the models, writing the same problem in different notation changes the embedding more than switching to any other problem in the benchmark. Giving the model the image together with the text helps little.

Our contributions are:
\begin{itemize}[leftmargin=*,itemsep=1pt,topsep=2pt]
\item a benchmark of 5,400 simulated trajectories, produced by one validated solver, each with a complete text description that never names the equation;
\item a reference distance with numerical error bounds, which show that every ranking we use is reliable, and prove it exactly for the linear equations;
\item an evaluation of three embedding models with text, image, and joint inputs, together with baselines and controls. Because only CPU hardware was available, the models were evaluated on 10 of the 100 samples per problem, chosen in advance.
\end{itemize}
The mathematical tools are classical: Euclidean distance, matrix exponentials, and standard time stepping. What is new is the controlled evidence.

\paragraph{Related work.}
APEBench provides configurable PDE benchmarks and reference solvers \citep{koehler2024apebench}. Physics-informed contrastive learning trains PDE representations with physical information \citep{lorsung2024picl}, and Unisolver and PDEformer condition learned solvers on the equation \citep{zhou2025unisolver,ye2024pdeformer}. We instead keep general-purpose encoders fixed and study the geometry they already have. Work on mathematical equivalence in embedding models, such as MELD, motivates our notation controls \citep{ye2026meld}. We extend that question from formulas to descriptions and images of evolving fields, where the ICs, BCs, and observation times all matter.
